\thispagestyle{myheadings}
\chapter{関連研究}
\label{sec:format}
% 本章ではインフォーマルなコミュニケーションの支援を行う研究と情報提示による行動変容支援に関する研究について本研究との関連を述べる．
本章では，人の位置情報を活用した行動分析および支援に関する研究を中心に，インフォーマルなコミュニケーションの支援，屋内における人の位置推定手法に関する研究，ならびに情報提示による行動変容支援について，本研究との関連を述べる．


\section{行動変容に関する研究}
本研究が目指すコミュニケーションの促進は，行動変容支援を目指す研究の議題の一種として位置づけられる．
行動変容に関する研究は，行動経済学におけるナッジ理論や，福祉・健康分野における生活行動支援など，多様な分野で展開されてきた．
近年では，IoT やセンシング技術の発展を背景に，人の行動や状況を把握した上で適切な介入を行う研究が増加している．
情報処理学会においてもIoT行動変容学研究グループが発足し，BTIW2023\cite{BTIW2023}という国際ワークショップが開催されるなど，IoTと組み合わせた行動変容に関する議論が進められている.


行動変容支援のアプローチは多岐にわたる．
例えば，行動経済学の知見に基づき選択肢の提示方法の工夫を通じて行動を誘導するナッジ理論に基づく研究\cite{nadge}\cite{nadge:body}が報告されている．
また，行動の結果を可視化・フィードバックとして提示し，自己認識を促して自発的な行動改善を支援する研究\cite{feedback:sit}も多い．
さらに，ポイント付与や達成目標の提示といったゲーミフィケーション要素と組み合わせ，行動への動機づけや継続性を高める試み\cite{game:health}も行われている．
福祉や健康支援の分野では，日常行動を継続的に記録し，ユーザの状況や状態に応じて適切なタイミングで介入を行い，生活習慣の改善や行動の定着を支援する研究が報告されている．
例えば，2型糖尿病の予防を目的として，活動量計とフラッシュグルコースモニターを組み合わせ，日常行動と血糖変動の関係を可視化して生活習慣の改善を促すデジタル行動変容介入が提案されている\cite{CGM}．
この研究では，行動の結果を生体情報として提示し，ユーザの自己認識を高めて自発的な行動調整を支援している．
これらの研究では，IoT やウェアラブルデバイスを用いた行動計測と組み合わせ，個人に適応した支援を実現する点が特徴である．


これらの研究に共通する点は，人の行動を直接的に制御するのではなく行動を選択するための情報や状況を整え，自発的な行動変化を促す点にある．
本研究における帰宅同行の促進も，特定の行動を強制するものではなく行動の選択に影響を与える情報を提示し，ユーザの行動変容を支援する試みとして位置づけられる．
% 本研究テーマでもこれらの手法を取り入れて課題の解決を目指す.

\section{インフォーマルなコミュニケーションを支援する研究}
\label{sec:format_thesis}
在室時のコミュニケーションの支援を行う研究が進んでいる．
例えば個別化が進むオフィスを対象としコミュニケーションのきっかけを提供するシステム\cite{cafe}がある．
これは個別スペースの訪問者に対してコーヒーを注ぎに行くという口実を与え，インフォーマルコミュニケーションを誘発している．
他にも共有インフォーマル空間を対象とした研究\cite{irori}がある．
この研究ではそこに行く理由や居るための理由となるオブジェクトに着目し，居心地の良い空間の創出からインフォーマルなコミュニケーションの発生を狙っている．


また，非対面でのコミュニケーションを支援する研究も行われている．
非対面環境では，対面時に自然に得られる他者の存在や活動状況に関する認識が失われやすいと指摘されており\cite{awareness}，こうした他者に関する認識が協調や円滑なコミュニケーションを支える重要な要素であるとされている.
非対面環境におけるコミュニケーション支援として，言語情報だけでなく非言語的な手がかりを用いて相手の存在や活動を伝える研究も行われている．
Go together\cite{gotogether}では遠隔にいる相手のジェスチャーや動作を簡略化して提示し，共在感を高める試みが報告されている．
一方で，こうした非対面での支援に加え，実際に同じ空間に集まる状況自体を促すアプローチも提案されている．


また，来訪促進を通してコミュニケーションの支援を行う研究がある．
在室履歴から未来の在室状況を予測する研究では現在の在室状況だけでなく未来の在室情報の予測により，次の訪問日時の検討を支援し来訪頻度の向上に繋げる．
曜日別の過去の在室の平均値と個人のスケジュールから未来の在室状況を予測する研究\cite{yohou}では，ユーザの推測が必要ない点で将来の在室可能性が数値で提示されている．
曜日別の過去の履歴を複数の行動パターンから形成されているとし，クラスタリングを通じて未来の在室状況を予測する研究\cite{hayashi}では，共通の話題や関心を持つグループに対して個人名を伏せて集まりやすい時刻を提示している．
これらの研究はいずれも，特定の個人，あるいは事前に定義されたグループを対象として来訪可能性を予測し，来訪行動そのものを支援する点に特徴がある．
一方，複数のユーザの行動を集合として扱い，グループ全体として行動が一致しやすいタイミングを導出するアプローチも考えられる．
ゲーミフィケーションと組み合わせて来訪を促進する研究\cite{ikusei}では，Login MATという育成アプリケーションの開発を通じてユーザの積極的な来訪の動機づけを行なっている．
支援を通じて対面でのコミュニケーション機会が増加するとコミュニティ全体の活性化に繋がると期待される．


このような在室時のコミュニケーション支援や来訪促進に関する研究は多く行われており，人が既に場に存在している状態を前提とした支援が中心である．
一方で，退室や帰宅といった行動転換のタイミングは在室中とは異なる心理的・社会的状況に位置づけられる可能性があるが，このタイミングを対象とした支援は必ずしも十分に検討されていない．
退室時の行動は業務外の時間として捉えられる場合が多く時間的にも限定された場面であるため，在室中の行動に比べて支援の対象として注目されにくかったためであると考えられる．



\section{情報提示による行動変容支援に関する研究}% モニターとかの画面提示以外の提示手法にも触れるべきかも，mizukichanのも環境の可視化による行動変容だしな
情報提示を通じて人の行動変容を支援する研究は，様々な形で行われている．
求める情報を投影するディスプレイを用いて対面での情報共有の促進を支援する研究\cite{huneas}\cite{concierge}では，共有スペースや待ち時間が発生する空間において，大型ディスプレイでの共通の興味のある話題の提示が，知らない人との会話数の増加に対して有効であると示されている．
これらの研究は，事前にユーザの興味情報の取得・設定を前提としている点で，本研究とは設計方針が異なる．


また，情報提示による行動変容支援はコミュニケーション促進に限らず，健康や生活行動の改善を目的とした研究にも広く用いられている．
例えば，ユーザの健康データや過ごす環境のデータの可視化により行動変容を狙った研究\cite{kenkou}\cite{kion}では，ユーザが意識していなかった情報の可視化により，ユーザの意識の変化に繋がっている．
これらの研究は，日常生活の中では把握しにくい身体状態や環境状態を可視化し，ユーザ自身の判断の支援を通じて自発的な行動選択の変化を促している．


さらに，情報提示にゲーミフィケーション要素を組み合わせる研究も報告されている．
例えば，心拍数などの身体情報をリアルタイムにグラフとして可視化し，既存のゲーム体験に反映させ娯楽性を保ちながら新たなゲーム性を付与する研究\cite{zinro}がある．
また，収集した温度や騒音などの環境情報をコレクション型のゲームとして提供し，ユーザに新たな環境への探索行動や身の回りの環境理解を促す研究\cite{kancolle}も行われている．
これらの研究は，情報を単に提示するだけでなく，遊びの要素を通じてユーザの内発的動機づけを高め，行動変容を持続的に支援する点に特徴がある．
以上より，情報提示はユーザの行動を直接制御するのではなく，可視化やゲーミフィケーションを通じて行動の判断や動機づけを支援し，自発的な行動変容を促す手法として位置付けられる．



\section{屋内で人の位置推定を行う手法に関する研究}
人の位置情報は様々な応用が期待されており，位置推定技術に対する需要は高い．
特に，人が日常生活や業務において屋内で過ごす時間は長く，屋内環境における位置情報の取得は重要な課題となっている．
近年，GPS を搭載したスマートフォンを利用した屋外位置推定が普及しているが，GPS は建物による遮蔽や反射の影響を受けやすく屋内環境において安定した位置の推定は困難である．
そのため，屋外位置推定における GPS に相当する，環境に依存せず広く適用可能である標準的な屋内位置推定手法が求められるが，未だ確立されていない．
こうした背景により，対象環境や利用目的に応じて様々なセンサや無線技術を用いた屋内位置推定手法が提案・研究されている．


屋内位置推定手法は，環境から直接位置を推定する絶対測位と，人の動作情報を用いて初期位置からの移動を推定する相対測位に分類され，それぞれに異なる利点と課題が存在する．
絶対測位の代表的な手法として，電波フィンガープリントを用いた位置推定や，複数の基準点からの距離情報を用いる三点測位方式が挙げられる．
電波フィンガープリントでは，事前に測位対象空間における Wi-Fi や BLE などの電波強度を含めた環境情報を収集し位置と対応させ，現在の測定値から位置を直接推定する\cite{fingerprint:WiFi}\cite{fingerprint:BLE}．
三点測位方式は位置が既知である複数の基地局やビーコンからの信号強度(RSSI)や到達時間をもとに距離を推定し，それらの交点として位置を算出する手法である\cite{threeposition:WiFi}\cite{threeposition:BLE}．
これらの絶対測位手法は，初期位置に依存せず位置推定が可能である反面，インフラの設置や環境変化への対応といった導入・運用コストが課題となる．
一方，相対測位の代表例として，加速度センサやジャイロセンサを用いた歩行者自律航法（PDR: Pedestrian Dead Reckoning）がある．
PDRは，初期位置を基準として歩行動作から移動量や方向を推定し，それらを順次足し合わせ現在位置を求める手法\cite{PDR:position}\cite{PDR:direction}である．
測位に特別なインフラを必要とせずスマートフォンなどの一般的な端末で利用可能である一方，推定誤差が時間とともに蓄積しやすいという課題がある．


屋内位置推定手法は測位精度の向上を図るほどインフラ設置や事前準備，運用管理にかかるコストが増大する傾向があり，精度とコストの間にはトレードオフが存在する．
しかし，屋内位置推定においては，利用目的や求められる測位精度に応じて，必ずしも高精度な位置推定が必要とは限らない．
例えば，在室管理や滞在状況の把握といった用途では位置を細かく推定する必要はなく，人物がどのエリアや部屋に存在するかを判別できれば十分な場合も多い．
このような目的に対して，BLE 信号強度に基づく近接判定（proximity）方式が提案，評価されている\cite{BLE:Proximity}．
位置精度は限定的であるものの，目的を限定する場合こうした導入や運用の負担が比較的小さい手法が選択される傾向にある．
他にも，単一の位置推定手法では精度や安定性に限界があるため，複数の手法を組み合わせて互いの欠点を補完するアプローチも多く用いられている．
例えば，相対測位であるPDRの誤差蓄積を抑制する手法として，外部情報を用いて推定位置を制約・補正するアプローチが用いられる場合が多い．
建物の構造を用いて物理的に不可能な経路とならないよう補正を行うマップマッチング\cite{mapmatching}や，既知の地点をランドマークとしてその地点で位置を再初期化するランドマーク測位\cite{landmark}が提案されている．
ただし，これらは事前準備や環境依存性といった新たなコストが発生する場合がある．
\section{人の位置情報を活用する研究}
人の位置情報は，行動分析や環境最適化など利用目的に応じて多様な分野で応用が検討されている．
特に屋内外における人の移動や滞在の履歴の取得により，人の行動特性や空間利用の実態を定量的に把握できる点が注目されている．
例えば，Google Maps\cite{GoogleMap}に代表されるナビゲーションシステムでは，位置情報を用いた現在地の可視化と経路案内が広く普及している．
また，位置情報とゲーミフィケーションを組み合わせた Pokémon GO\cite{PokemonGo}のようなゲームでは，実空間での移動をゲーム体験と結び付けユーザの行動変容を促す仕組みが実用化されている．
これらは人の移動や滞在といった行動を位置情報から捉え，その特徴を分析・活用している点に共通点がある．
こういった，人の行動特性を明らかにする行動分析に関する研究が数多く行われている．


行動分析の分野では，位置履歴を用いてエリアごとの滞在時間や訪問頻度を分析し，人の行動傾向や利用パターンを明らかにする研究が数多く行われている．
大規模商業施設において来訪者の歩行軌跡を分析し，空間の実際の利用状況と設計意図との差異を示した研究\cite{layout}が報告されている．
さらに，長期間の位置情報の蓄積から個人や集団の行動パターンを抽出し，行動の類型化や特徴量の分析を行う試みもある．
例えば，人物動線データ群から行動パターンを分類し，通常行動から逸脱した行動を検出する研究\cite{pattern}が報告されている．
また，大地震時を想定した広域避難シミュレーションにおいて，避難軌跡をクラスタリングして行動傾向を把握し，将来の避難行動予測へ応用する研究\cite{simulation}も提案されている．
このように，人の位置情報を時系列に蓄積・分析すると，日常行動から非常時の避難行動に至るまで様々な状況における行動傾向や空間利用の特徴を明らかにできると示されており，位置情報に基づく行動分析は幅広い分野で有効である．


また，人の位置情報は利用目的や環境に応じて多様な形で活用可能であり，本研究で前提となる部屋単位の滞在管理もこうした屋内位置推定の技術における応用の一例として位置づけられる．
滞在管理は，人物がいつどの空間に存在しているかの把握を目的としており，詳細な位置や移動経路を高精度に推定する必要がない点に特徴がある．
特に，オフィスや研究室，大学施設などの屋内環境では，在室状況の把握が円滑な業務遂行やコミュニケーションの促進に寄与すると考えられており，前述した来訪促進を目的とする研究もこの一例である．
この情報は在室人数の把握による空間利用状況の可視化，利用頻度や滞在時間の分析に基づく運用改善，在室情報の共有による対面コミュニケーションの支援などを通じて活用される．


さらに，こうして得られた人の位置や滞在状況の情報は，環境最適化を目的とした研究にも応用されている．
位置情報を用いて空調，照明，案内表示などの制御を行うシステムが提案されており，ユーザの存在や分布に応じて環境を動的に調整する試みが行われている．
例えば，滞在状況に応じて照明を制御する知的照明システム\cite{light}や，人の分布に基づいて空調を制御する手法\cite{airControl}などが検討されている．
これらは快適性の向上やエネルギー消費の削減を目的としている．
このように，人の位置情報は分析，最適化，可視化，娯楽など，多様な目的で活用可能な情報であり，屋内外を問わず幅広い応用が検討されている．



% 卒業論文の作成には \LaTeX やWordなどのソフトウェアの使用を原則とする．
% 通常，論文は2部作成し，1部は情報科学部事務室に提出し，1部は各自が保管
% する．事務室に提出する論文は，A4サイズファイルに綴じ，指導教員に確認の
% 印を押してもらう．綴じ込み用フラットファイルは，1論文につき1冊，学科か
% ら各教員へ支給される．ファイルの表紙には論文題目と学籍番号・氏名，指導
% 教員名を明記する．ファイルを開いたときの第1頁目にも，ファイルの表紙と
% 同一項目を明記した中表紙をつける．これは，ファイルの外表紙が破損や退色
% をしても，内容がわかるようにするためである．

% \begin{tabular}{llcl}
%   (1) & 提出部数                & : & 1部(学部事務室に提出)                         \\
%   (2) & 用紙                  & : & A4サイズ白紙(罫線・枠等なし)                     \\
%   (3) & 印字書式                &   &                                      \\
%       & \hspace*{1zw}・印字方向  & : & 横書き                                  \\
%       & \hspace*{1zw}・段組み数  & : & 1段                                   \\
%       & \hspace*{1zw}・文字サイズ & : & 10.5ポイント程度                           \\
%       & \hspace*{1zw}・1頁の行数 & : & 45行程度                                \\
%       & \hspace*{1zw}・頁番号印刷 & : & する(ページ外側の上部)                         \\
%       & \hspace*{1zw}・上余白   & : & 26mm程度                               \\
%       & \hspace*{1zw}・下余白   & : & 26mm程度                               \\
%       & \hspace*{1zw}・左余白   & : & 30mm程度                               \\
%       & \hspace*{1zw}・右余白   & : & 20mm程度                               \\
%   (4) & 分量                  & : & 制限なし                                 \\
%   (5) & 図表                  & : & ・図，表にはそれぞれ通し番号とタイトルをつける              \\
%       &                     &   & ・\underline{表番号およびタイトルは表の上側につける}     \\
%       &                     &   & ・\underline{図番号およびタイトルは図の下側につける}     \\
%       &                     &   & ・グラフには単位および軸の意味を記し，見て理解できるよ
%   うにする                                                                 \\
%       &                     &   & ・図，表は必ず本文から参照し説明する                   \\
%       &                     &   & \hspace*{1zw}(図，表だけで理解できるのが望ましい)     \\
%   (6) & 見出し                 & : & ・章，節の番号表記，見出しの書き方は別紙に付けたサ
%   ンプルに従う                                                               \\
%       &                     &   & ・章が変わるところでは改ページをおこなう                 \\
%       &                     &   & ・節が変わるところでは見出しの前に1行あける               \\
%   (7) & 文字種の使い分け            & : & ・文章中の仮名および漢字は全角とする                   \\
%       &                     &   & ・数字，アルファベットは半角とする                    \\
%       &                     &   & ・章，節などの見出しは太字などでわかりやすくする             \\
%       &                     &   & ・網掛け，白抜き，倍角等，冗長な強調表現は避ける             \\
%   (8) & 参考文献                & : & ・参考文献の書き方は付録に示すが，指導教員の
%   指示がある場合は，                                                            \\
%       &                     &   & \hspace*{1zw}それに従うこと                 \\
%   (9) & タイトル等               & : & ・研究内容を適切に表すタイトルをつける．                 \\
%       &                     &   & \hspace*{1zw}(2行を越えるなど長すぎないように気をつける) \\
%       &                     &   & ・必要に応じて副題を付けても良いが，この場合も長さに気
%   をつける                                                                 \\
%       &                     &   & ・綴じ込み用ファイルにつける外表紙と，中表紙の用意する          \\
%       &                     &   & \hspace*{1zw}(同じ書式で構わない)             \\
%       &                     &   & ・指導教員から合格を受けたら，中表紙の教員名の右横に印
%   鑑をもらう                                                                \\
%       &                     &   & ・表紙のサンプルを別紙に示す                       \\
%       &                     &   & ・目次をつける                              \\
%       &                     &   & ・複数人で書いた場合は執筆範囲が担当がわかるようにする          \\
% \end{tabular}

% \section{論文の構成例}

% 卒業論文の構成については各指導教員の指示に従うことになるが，参考までに
% 構成例を次に示す．

% \begin{description}
%   \item[第1章] 「はじめに」または「序論」として，論文概要及び論文の構成
%         について説明する．
%   \item[第2章] 「背景」として論文の背景にあたる内容を書く．
%   \item[第3章] 「提案手法」として，\textbf{研究の目的}からアプローチまで
%         を説明する．章のタイトルが，卒論のタイトルと一致する場合もある．
%   \item[第4章] 「実験及び考察」提案手法が研究の目的を達成できているかど
%         うかを評価確認し，考察する．
%   \item[第5章] 「まとめと今後の課題」として，達成できたことと，今後の課
%         題として取り組むべき内容についてまとめる．
%   \item[謝辞] お世話になった先生方や先輩達へのお礼を述べる．
%   \item[参考文献] 卒業論文の執筆に際し，参考にした文献について記述する．
%         本文中と\textbf{相互参照}する．
%   \item[付録] 本文中に入れることが困難であった詳細な定理証明や，アルゴリ
%         ズムなどを記述する．
% \end{description}
% これら，あくまで例であり，各研究室の指導教員の指示に従う．

% \section{論文の製本形態と関連資料の整理}
% \label{sec:style}

% 指導教員へ提出する論文の提出形態は指導教員の指示に従う．事務室へ提出す
% るものと同等のものを必要とされる場合もあるが，電子的に成果物を提出する
% 場合もある．また本論文以外のデータやプログラム，作品等についても研究室
% に残すものとし，不正等が認められる場合は，厳正な処分が下されることもあ
% り得る．

% \section{レジュメ(概要)の書式}
% \label{sec:format_abst}

% レジュメの \LaTeX やWordなどのテンプレートは配布されたものを用いて良い．
% また，レジュメは電子的にPDFで提出するため，ソフトウェアの種類を限定す
% るものではない．

% なお，レジュメはテンプレートで示される書式に従うこととするが，伝統的に
% は次のような書式が用いられている．ここで示す書式は，一応の目安であり，
% 1行の文字数や段間の空白については変更が許される．ソフトウェアによって
% は，1行の文字数が固定できないものがある．その場合は，ほぼこの書式に準
% ずる文字数になるよう，フォントサイズや文字間隔，行間隔を工夫する．

% \begin{tabular}{llcl}
%   (1) & 提出部数                    & : & 1部(学部事務室に提出)                         \\
%   (2) & 用紙                      & : & A4サイズ白紙(罫線・枠等なし)                     \\
%   (3) & 印字書式                    &   &                                      \\
%       & \hspace*{1zw}・印字方向      & : & 横書き                                  \\
%       & \hspace*{1zw}・段組み数      & : & 2段                                   \\
%       & \hspace*{1zw}・表題文字サイズ   & : & 16ポイント程度                             \\
%       & \hspace*{1zw}・節見出し文字サイズ & : & 12ポイント程度                             \\
%       & \hspace*{1zw}・本文文字サイズ   & : & 10ポイント程度                             \\
%       & \hspace*{1zw}・1頁の行数     & : & 45行程度                                \\
%       & \hspace*{1zw}・頁番号印刷     & : & しない                                  \\
%       & \hspace*{1zw}・上余白       & : & 15mm程度                               \\
%       & \hspace*{1zw}・下余白       & : & 20mm程度                               \\
%       & \hspace*{1zw}・左余白       & : & 15mm程度                               \\
%       & \hspace*{1zw}・右余白       & : & 15mm程度                               \\
%   (4) & 分量                      & : & 2頁                                   \\
%   (5) & 図表                      & : & ・論文の書式に準ずるものとする                      \\
%   (6) & 見出し                     & : & ・論文の書式に準ずるが，章が変わるところでは，大見出し
%   の前に                                                                      \\
%       &                         &   & \hspace*{1zw}1行あけるものとする              \\
%   (7) & 文字種の使い分け                & : & ・論文の書式に準ずるものとする                      \\
%   (8) & 参考文献                    & : & ・論文の書式に準ずるものとする                      \\
%   (9) & タイトル等                   & : & ・研究内容を適切に表すタイトルをつける．                 \\
%       &                         &   & \hspace*{1zw}(2行を越えるなど長すぎないように気をつける) \\
%       &                         &   & ・必要に応じて副題を付けても良いが，この場合も長さに気
%   をつける                                                                     \\
%       &                         &   & ・タイトルに続いて学籍番号，名前をつける                 \\
%       &                         &   & ・最後に指導教員名をつける                        \\
% \end{tabular}

% \clearpage

% Local Variables: 
% mode: japanese-LaTeX
% TeX-master: "root"
% End: 
